\section{Main Theorem}

\begin{theorem}[Endpoint conditional anti-concentration]
\label{thm:main}
Under Assumptions~\ref{assump:compact-parameter-domain}--
\ref{assump:mean-endpoint-conditional-caps}, fix the finite
$\eta=(\bar\kappa_0,\bar\kappa_\infty)$.  For every $d\geq1$, $R\geq1$,
$\mu\in\mathcal D_{d,R,\eta}$, and $I\in\mathcal I(\Theta)$,
\[
\begin{aligned}
\mu(H_{d,I})
&\leq \bar\kappa_0B_0(d,R)|I_0|
 +\bar\kappa_\infty B_\infty(d,R)(|I_+|+|I_-|)\\
&\leq M_\eta(d,R)|I|.
\end{aligned}
\]
Consequently,
\[
\boxed{
C_{\mathcal D_{d,R,\eta}}
\leq M_\eta(d,R)
\leq P_\eta(d,R)
=\bar\kappa_*d+\frac{\bar\kappa_*}{2}Rd^2.}
\]
The middle quantity is the maximum of the two chart coefficients, rather
than their sum.

Separately, if $\bar\kappa_0,\bar\kappa_\infty\geq1/2$, then for every
$d\geq1$ and $R\geq1$ the law $\mu^{\mathrm{wit}}_{d,R}$ belongs to
$\mathcal D_{d,R,\eta}$, and
\[
K_0^{\mu^{\mathrm{wit}}_{d,R}}
=K_\infty^{\mu^{\mathrm{wit}}_{d,R}}
=\frac1{2R}
\quad\text{almost surely}.
\]
This threshold is not required for either root-hitting inequality.

All exposed dependence is on $d$, $R$, $\bar\kappa_0$,
$\bar\kappa_\infty$, and the displayed interval lengths; there are no
hidden constants.  The pair $\eta$ and the domain $\Theta$ are fixed, and
the bound itself is independent of $\Theta$.  The statement is a
deterministic inequality for each probability $\mu(H_{d,I})$, followed by
deterministic suprema, with no confidence parameter.  It is static and
uniform over all positive-length intervals, with no horizon or limiting
qualification.  Its norm and metric inputs are the conditional
$L^\infty(\R)$ density norms and one-dimensional Lebesgue interval length.
\end{theorem}

\begin{corollary}[Fixed-cap polynomial rate]
\label{cor:fixed-cap-rate}
Under Assumptions~\ref{assump:compact-parameter-domain}--
\ref{assump:mean-endpoint-conditional-caps}, for the single fixed finite
$\eta$ and every $d\geq1$, $R\geq1$,
\[
C_{\mathcal D_{d,R,\eta}}
\leq \bar\kappa_*d+\frac{\bar\kappa_*}{2}Rd^2.
\]
This is a pointwise, deterministic, static bound with no hidden constant or
probability conversion; its displayed polynomial has total degree at most
three in $(d,R)$ and coefficients depending only on $\eta$.
\end{corollary}

\begin{proof}
Theorem~\ref{thm:main} gives
$C_{\mathcal D_{d,R,\eta}}\leq M_\eta(d,R)$.  The complete scalar
specialization, including the choices
$\bar\kappa_*=\max\{\bar\kappa_0,\bar\kappa_\infty\}$, the conditions
$d\geq1$ and $R\geq1$, the inequalities $1\leq d$ and
$d(d-1)\leq d^2$, and the absence of any probability-mode conversion, is
proved in Proposition~\ref{prop:fixed-eta-polynomial-specialization}.
Combining the two statements gives the displayed rate.
\end{proof}
